\section{A boundary connectivity lemma}\label{s:topology}

In this appendix we prove the lattice-topology fact used in \S\ref{s:cutgeom}. Two sites $x, x' \in \Z^2$ are \emph{neighbours} if $|x - x'|_1 = 1$ and \emph{$*$-neighbours} if $|x - x'|_\infty = 1$. A set is \emph{connected} (respectively \emph{$*$-connected}) if any two of its points are joined by a path of neighbours (of $*$-neighbours) inside the set. A \emph{block} is a set $p + \{0, 1\}^2$, $p \in \Z^2$; two sites of a block that are not neighbours are called \emph{diagonally opposite}. For a finite set $A \subset \Z^2$ let
\begin{equation*}
\partial_\infty A = \{ x \in \Z^2 \setminus A : \operatorname{dist}_\infty(x, A) = 1 \} .
\end{equation*}

\begin{lemma}\label{l:boundaryconn}
Let $A \subset \Z^2$ be finite and nonempty, and assume that
\begin{enumerate}
\item[(i)] $A$ is $*$-connected;
\item[(ii)] $\Z^2 \setminus A$ is $*$-connected;
\item[(iii)] no block meets $A$ in exactly two diagonally opposite sites.
\end{enumerate}
Then $\partial_\infty A$ is connected.
\end{lemma}

Without (iii) the conclusion fails: for $A = \{0, 1, 2\}^2 \setminus \{(0, 0), (1, 1)\}$, (i) and (ii) hold, but $(1, 1) \in \partial_\infty A$ has no neighbour outside $A$, so $\partial_\infty A$ is not connected.

\begin{proof}
Let $\mathsf K$ be the graph on $\Z^2$ whose edges join $*$-neighbours. An \emph{edge set} is a finite set of edges of $\mathsf K$; edge sets form a vector space over $\mathbb F_2$ under symmetric difference, which we write as a sum. An edge set is \emph{even} if every site lies in an even number of its edges. A \emph{triangle} is a set of three pairwise $*$-neighbouring sites; its first coordinates lie in an interval of length $1$, and so do its second coordinates, so a triangle consists of three sites of one block. For a triangle $T$ let $\partial T$ be the set of its three edges, and for a block $P$ let $\partial P$ be the set of the four edges of $\mathsf K$ that join neighbours in $P$.

\emph{Step 1: every even edge set is a sum of sets $\partial T$.} Let $Z$ be even. For every edge $e = \{p, p + (1, \epsilon)\}$ of $Z$ with $\epsilon \in \{1, -1\}$ (a diagonal edge), let $T_e = \{p, p + (1, 0), p + (1, \epsilon)\}$. The only diagonal edge of $\partial T_e$ is $e$, so $Z' = Z + \sum_e \partial T_e$ is an even edge set without diagonal edges. For $(i, j) \in \Z^2$ let $P_{ij} = (i, j) + \{0, 1\}^2$, and let $n(i, j) \in \mathbb F_2$ be the parity of the number of $j' > j$ such that $\{(i, j'), (i + 1, j')\} \in Z'$. The number of edges of $Z'$ of the form $\{(i, j'), (i + 1, j')\}$, $j' \in \Z$, is even, because it has the parity of the sum over the sites $x$ with $x_1 \le i$ of the number of edges of $Z'$ at $x$. Hence $n(i, j) = 0$ for all but finitely many $(i, j)$. We claim that $Z' = \sum_{i, j} n(i, j) \, \partial P_{ij}$. The horizontal edge $\{(i, j), (i + 1, j)\}$ lies in $\partial P_{ij}$ and $\partial P_{i, j - 1}$ only, and $n(i, j - 1) + n(i, j)$ is $1$ if and only if this edge is in $Z'$. The vertical edge $\{(i, j), (i, j + 1)\}$ lies in $\partial P_{i - 1, j}$ and $\partial P_{ij}$ only, and $n(i - 1, j) + n(i, j)$ is the parity of the number of horizontal edges of $Z'$ at the sites $(i, j')$, $j' \ge j + 1$. Since $Z'$ is even, this is also the parity of the number of vertical edges of $Z'$ at these sites, counted with multiplicity; every edge $\{(i, j'), (i, j' + 1)\}$ with $j' \ge j + 1$ is counted twice and $\{(i, j), (i, j + 1)\}$ once, which proves the claim. Finally, $\partial P_{ij} = \partial T + \partial T'$ for the triangles $T = \{(i, j), (i + 1, j), (i + 1, j + 1)\}$ and $T' = \{(i, j), (i, j + 1), (i + 1, j + 1)\}$.

\emph{Step 2: the boundary edges are linked.} Let $\mathcal B$ be the set of edges of $\mathsf K$ with one endpoint in $A$ and the other outside $A$; it is nonempty. Call two edges of $\mathcal B$ \emph{linked} if they lie in $\partial T$ for a common triangle $T$. We claim that $\mathcal B$ is connected for this relation. Suppose instead that $\mathcal B = \mathcal B_1 \cup \mathcal B_2$ with $\mathcal B_1, \mathcal B_2$ disjoint and nonempty and no edge of $\mathcal B_1$ linked to an edge of $\mathcal B_2$. For an edge set $Z$ let $\varphi(Z) \in \mathbb F_2$ be the parity of $|Z \cap \mathcal B_1|$; $\varphi$ is linear. For every even $Z$ the number $|Z \cap \mathcal B|$ is even, since it has the parity of the sum over $x \in A$ of the number of edges of $Z$ at $x$. For a triangle $T$, $\partial T \cap \mathcal B$ lies in $\mathcal B_1$ or in $\mathcal B_2$, so $\varphi(\partial T)$ is the parity of $|\partial T \cap \mathcal B|$ or $0$, and in both cases $\varphi(\partial T) = 0$. By Step 1, $\varphi(Z) = 0$ for every even $Z$. Now choose $\{a_1, b_1\} \in \mathcal B_1$ and $\{a_2, b_2\} \in \mathcal B_2$ with $a_1, a_2 \in A$. By (i) there is a path of $*$-neighbours from $a_1$ to $a_2$ in $A$, and by (ii) one from $b_2$ to $b_1$ outside $A$. Together with the edges $\{a_2, b_2\}$ and $\{b_1, a_1\}$ they form a closed path, and the set $Z$ of edges that it uses an odd number of times is even. The two paths use no edge of $\mathcal B$, so $Z \cap \mathcal B = \{\{a_1, b_1\}, \{a_2, b_2\}\}$ and $\varphi(Z) = 1$, a contradiction.

\emph{Step 3: linked edges have connected outer endpoints.} Every edge of $\mathcal B$ has exactly one endpoint outside $A$, its \emph{outer endpoint}, which lies in $\partial_\infty A$, and every site of $\partial_\infty A$ is the outer endpoint of an edge of $\mathcal B$. Let $\{a, b\}, \{a', b'\} \in \mathcal B$ be linked, with $a, a' \in A$, through a triangle $T$ contained in a block $Q$. We claim that $b$ and $b'$ are equal, or neighbours, or have a common neighbour in $\partial_\infty A$. Suppose that $b \neq b'$ are not neighbours. Since $T$ has three sites and contains $a$, $b$, $b'$, we have $T = \{a, b, b'\}$, $a' = a$, and $b$, $b'$ are diagonally opposite in $Q$. Let $d$ be the fourth site of $Q$; then $a$ and $d$ are diagonally opposite. Since $b, b' \notin A$, (iii) gives $d \notin A$. As $d$ is a $*$-neighbour of $a \in A$, we have $d \in \partial_\infty A$, and $d$ is a neighbour of both $b$ and $b'$.

Let $x, x' \in \partial_\infty A$ be the outer endpoints of $e, e' \in \mathcal B$. By Step 2 there are $e = e_0, e_1, \dots, e_n = e'$ in $\mathcal B$ with $e_{m - 1}$ and $e_m$ linked, and by Step 3 the outer endpoints of $e_{m - 1}$ and $e_m$ are joined by a path of neighbours in $\partial_\infty A$ of length at most $2$. Hence $x$ and $x'$ are joined by a path of neighbours in $\partial_\infty A$.
\end{proof}

The parity argument of Step~2 is the one of Tim\'ar's proof of boundary connectivity \cite{Tim13}; we included it to keep the statement in the form needed here.
